\documentclass[10pt,a4paper]{amsart}
\usepackage[margin=19mm]{geometry}
\usepackage{amsmath,amssymb,mathtools}
\usepackage[scr=rsfs,cal=euler]{mathalfa}
\usepackage{xfrac}
\usepackage[unicode,hidelinks,bookmarks=false]{hyperref}
\newcommand{\ie}{i.e.\ }
\newcommand{\cf}{cf.\ }
\newcommand{\resp}{respectively\ }
\newcommand{\Subsection}{\subsection}
\providecommand{\nobreakdash}{\nobreak\hbox{-}}
\setlength{\emergencystretch}{2em}
\multlinegap0pt
\usepackage{bm}
\usepackage{bbm}
\usepackage{dsfont}
\usepackage{xspace}
\usepackage{cancel}
\usepackage{mathtools}
\usepackage{cleveref}
\usepackage{amsthm}
\makeatletter
\@ifundefined{theorem}{\newtheorem{theorem}{Theorem}}{}
\@ifundefined{definition}{\newtheorem{definition}[theorem]{Definition}}{}
\@ifundefined{proposition}{\newtheorem{proposition}[theorem]{Proposition}}{}
\makeatother
\newcommand{\N}{\mathbb{N}}
\newcommand{\R}{\mathbb{R}}
\newcommand{\SE}{\mathbb{S}}
\newcommand{\Z}{\mathbb{Z}}
\newcommand{\lowerent}{\underline{\operatorname{ent}}}
\newcommand{\spa}{\operatorname{spa}}
\newcommand{\upperent}{\overline{\operatorname{ent}}}
\title[Topological slow entropy, sequence entropy, and generalized ]{Topological slow entropy, sequence entropy, and generalized $[T,T^{-1}]$ systems}
\author{Nicanor Carrasco-Vargas}
\date{Source: arXiv:2506.17932v1 (CC BY 4.0)}
\begin{document}
\maketitle

\noindent\emph{Source and licence.} Topological slow entropy, sequence entropy, and generalized $[T,T^{-1}]$ systems. Version arXiv:2506.17932v1 (CC BY 4.0), distributed under the Creative Commons Attribution 4.0 International licence, as recorded in the arXiv licence metadata for this exact version and shown on the abstract page https://arxiv.org/abs/2506.17932v1. Full rights notice: licenses/arxiv-2506.17932-CC-BY-4.0.txt.\par\smallskip

\noindent\textit{Editorial scope.} Counted unit: the Proposition whose source label is \texttt{prop:entropy} in the article (source0.tex lines 649--652, source label \texttt{prop:entropy}), delivered together with the article's own complete original proof of it (source0.tex lines 653--712, one \texttt{proof} environment of 60 lines). No mathematics is written, repaired or re-derived: statement and proof are the article's own text line for line. Required context retained uncounted: thm:entropy (lines 111--114); monotony-properties-of-spa (lines 206--212); constant-C\_m (lines 429--431); estimate (lines 432--440); constant-m (lines 545--545); def:good (lines 564--570); def:scales (lines 578--587); prop:the-set-A (lines 606--610). Every definition, display and label of those lines is retained unchanged. Omitted from this package and not redistributed: 1--110, 115--205, 213--428, 441--544, 546--563, 571--577, 588--605, 611--648, 713--881. Nothing inside a retained excerpt is omitted or altered: every retained body is delivered exactly as the article's lines are written, so no omission receipt of any kind accompanies this package. Citations: the retained excerpts carry no citation command at all, so no bibliography entry of the article is transcribed and no reference is redistributed. Reference adaptation: none was needed and none was made; every reference inside the delivered unit resolves inside it, so no unresolved reference or citation is produced. Printed slips kept literal: no printed word, symbol, hypothesis or number of the counted statement or of its proof was altered. Layout and class: the article's own class is \texttt{amsart} with packages including comment, geometry, amsmath, flexisym, breqn, amsfonts, amssymb, graphicx, hyperref, todonotes, cleveref; the wrapper is the lane's pinned \texttt{amsart} editorial wrapper at 10pt on A4, which restates the theorem-like environments the retained text needs and adds the article's own macro definitions that the retained text needs (\texttt{\textbackslash N}, \texttt{\textbackslash R}, \texttt{\textbackslash SE}, \texttt{\textbackslash Z}, \texttt{\textbackslash lowerent}, \texttt{\textbackslash spa}, \texttt{\textbackslash upperent}), with extra wrapper packages (bm, bbm, dsfont, xspace, cancel, mathtools, cleveref). The delivered numbering is the unit's own. Assets: no retained span contains a figure environment, a graphics-inclusion command, a PDF-page inclusion, a table or a TikZ picture; no image file is copied into this package, the package declares no asset dependency and no image file is redistributed with it. Licence: version arXiv:2506.17932v1 of this article is distributed under the Creative Commons Attribution 4.0 International licence, as recorded in the arXiv licence metadata for this exact version and shown on the abstract page https://arxiv.org/abs/2506.17932v1. A full rights notice is delivered at licenses/arxiv-2506.17932-CC-BY-4.0.txt. Characters: every retained excerpt is pure ASCII, so the delivered engine, XeLaTeX with its default Latin Modern text font, reports no missing character.
\begin{theorem}\label{thm:entropy}
Let $(Y,S)$ be a subshift, and let $\tau\colon Y\to \Z$ be a continuous map which is $\lambda$-unbounded for some $\lambda>0$. Then there exists a scale $\textbf{a}=\{a_n(t)\}_{n\in\N,t>0}$ depending only on $(Y,S,\tau)$, and such that
\[\upperent_{\textbf{a}}(S\rtimes_\tau T)=\lowerent_{\textbf{a}}(S\rtimes_\tau T)=h_{top}(T)\]  for every invertible topological dynamical system $(X,T)$.
\end{theorem} 

\begin{proposition}\label{monotony-properties-of-spa}Let $F,F'\Subset\SE$ and $n\in\SE$. 
    \begin{enumerate}
        \item If $F\subset F'$ then $\spa(T,F,\epsilon)\leq\spa(T,F',\epsilon)$. 
        \item We have $\spa(T,F,\epsilon)=\spa(T,F+n,\epsilon)$. 
        \item If $F$ and $F'$ are intervals and $|F|\leq |F'|$ then $\spa(T,F,\epsilon)\leq \spa(T,F',\epsilon)$
    \end{enumerate}
    \end{proposition}

\begin{definition}\label{constant-C_m} Let $m\in\N$. A set $F\Subset \Z$ is said to have $m$-\textbf{bounded gaps} if the set $F+\{0,\dots,m-1\}$ is an interval. We define $C_m(F)\in\R $ by $
C_m(F)=|F+\{0,\dots,m-1\}|/|F|$.
\end{definition} The main result of this subsection is the following.

\begin{proposition}\label{estimate}
Fix $m\in\N$. Suppose that $h_{top}(T)<\infty$. Then for every  $\epsilon>0$ we can choose $\delta\in (0,\epsilon)$ and $n_0\in\N$ such that
\begin{enumerate}
    \item $\frac{1}{|F|}\log \spa(T,F,\delta)\leq C_m(F) (h_{top}(T)+\epsilon)$ for every $F\Subset\Z$ with $m$-bounded gaps and $|F|\geq n_0$,
    \item $
    \frac{1}{|F|}\log \spa(T,F,\delta)\geq C_m(F) (h_{top}(T)-\epsilon)$ for every $F\Subset\Z$.
\end{enumerate}
    If $h_{top}(T)=\infty$ then for every $N\in\N$ we can choose $\delta>0$ such that $C_m(F)\cdot N\leq \frac{1}{|F|}\log \spa(T,F,\delta)$ for every $F\Subset\Z$.
\end{proposition} 

\begin{equation}\label{constant-m}|\tau(y)|\leq m \ \text{ for all } y\in Y.\end{equation}

\begin{definition}\label{def:good}\label{lambda-unbounded}
    Given $w\in L_{n,s}(Y)$ we define $r_n(w)$ by 
    \[r_n(w)=|\tau^{\{0,\dots,n-1\}}(w)|.\]
    Given $\lambda>0$, we say that $\tau$ is $\lambda$-unbounded if for
    all $N\in\N$ we have
\[\liminf_{n\in\N}\frac{|\{w\in L_{n,s}(Y) : r_n(w)\geq N\}|}{|L_{n,s}(Y)|}\geq \lambda.\]
\end{definition}

\begin{definition}\label{def:scales}
Let $w\in L_{n,s}(Y)$, and recall from \Cref{def:good} that
\[r_n(w)= |\tau^{\{0,\dots,n-1\}}(w)|.\]
We define $q_n(w)$ by 
\[q_n(w)=r_n(w)\cdot C_m(\tau^{\{0,\dots,n-1\}}(w))\]
Here $C_m$ is the constant defined in \Cref{constant-C_m}, and $m$ was chosen in \Cref{constant-m}. We define the scale $a_n(t)$ by
\[a_n(t)=\sum_{w\in L_{n,s}(Y)}e^{t\cdot q_n(w)}.\]
Furthermore, given an arbitrary scale $b_n(t)$, we define the scale $c_n(t)$ by 
\[c_n(t)=\sum_{w\in L_{n,s}(Y)}b_{r_n(w)}(t).\]
\end{definition}

\begin{proposition}\label{prop:the-set-A}
    Let $\epsilon\in (0,2^{-s-1})$ and $n\geq 1$. Then we have the inequalities
    \[A_n(2\epsilon)\leq \spa(S\rtimes_\tau T,\{0,\dots,n-1\},\epsilon)\leq E_{\epsilon}\cdot  A_n(\epsilon/2),\]
    where $E_{\epsilon}$ is a constant depending on $\epsilon$ and not on $n$.
\end{proposition}

\begin{proposition}[Proof of \Cref{thm:entropy}]\label{prop:entropy}
The upper and lower slow entropy of $S\rtimes_\tau T$ with respect to $a_n(t)$ is equal to $h_{top}(T)$. That is,
\[\lowerent_{\textbf{a}}(S\rtimes_\tau T)=\upperent_{\textbf{a}}(S\rtimes_\tau T)=h_{top}(T)\]
\end{proposition}

\begin{proof}
    We start by proving the inequality $
    \upperent_{\textbf{a}}(S\rtimes_\tau T)\leq h_{top}(T)$. This is clearly true if $h_{top}(T)=\infty$, so we assume $h_{top}(T)<\infty$. Let $t>0$ be arbitrary with $h_{top}(T)<t$. We will prove that $\upperent_{\textbf{a}}(S\rtimes_\tau T)\leq t$. According to the definition of upper slow entropy, we must show that for arbitrarily small $\epsilon$ we have
    \[\limsup_{n\to\infty} \frac{\spa(S\rtimes_\tau T,\{0,\dots,n-1\},\epsilon)}{a_n(t)}=0.\]
    By \Cref{prop:the-set-A} it suffices to prove that for arbitrarily small $\epsilon$
    \begin{equation}\label{eq:limsup}\limsup_{n\to\infty} \frac{A_n(\epsilon)}{a_n(t)}=0.\end{equation}
    Choose $t'$ with $h_{top}(T)<t'<t$. \Cref{estimate} shows that there is an arbitrarily small $\epsilon$, fixed from now on, with the following property. There is $n_0\in\N$ such that \[\spa(T,F,\epsilon)\leq e^{C_m(F)|F|t'}\]
    for every $F\Subset\Z$ with $m$-bounded gaps and $|F|\geq n_0$. Recall that $r_n(w)$ and $q_n(w)$ are defined in \Cref{def:scales}. If we apply the inequality above to $F=\tau^{\{0,\dots,n-1\}}(w)$ for some $w\in L_{n,s}(Y)$ with $r_n(w)\geq n_0$, we obtain
    \begin{equation}\label{eq:estimate-1} \spa(T,\tau^{\{0,\dots,n-1\}}(w),\epsilon)\leq e^{q_n(w)t'}.\end{equation}
    We will prove that for this $\epsilon$ the limsup in \Cref{eq:limsup} is 0. For this purpose we fix an arbitrary $\gamma>0$ and we prove that the limsup is at most a multiple of $\gamma$. We will now choose some parameters, the reason of these choices will become clear later. 
    
    Since $t-t'>0$, the function $1/e^{n(t-t')}$ converges to $0$ as $n\to\infty$, and we can find $n_1\geq n_0$ such that for all $n\geq n_1$ we have $1/e^{n(t-t')}<\gamma$. For a word $w\in L_{n,s}(Y)$ with $r_n(w)\geq n_1$, by \Cref{eq:estimate-1} we have \[\spa(T,\tau^{\{0,\dots,n-1\}}(w),\epsilon)/e^{q_n(w)t}\leq 1/{e^{q_n(w)(t-t')}}<\gamma.\] The relevant consequence for us is 
    \begin{equation}\label{n1-definition}
    \frac{\spa(T,\tau^{\{0,\dots,n-1\}}(w),\epsilon)}{e^{q_n(w)t}}<\gamma, \ \ \ r_n(w)\geq n_1.
    \end{equation}
    After choosing $n_1$, we let $N\in\N$ be equal to $\spa(T,I,\epsilon)$, where $I\subset\Z$ is an interval with length $2n_1m$. Therefore when  $F\Subset\Z$ is an arbitrary set with $m$-bounded gaps and $|F|\leq n_1$, a translate of $F$ is contained in $I$ and therefore $\spa(T,F,\epsilon)\leq N$ (\Cref{monotony-properties-of-spa}).  
    
    After choosing $N$ we observe that the function $N/e^{nt}$ converges to $0$ as $n\to\infty$, and then we can find $n_2\geq n_1$ with the property 
    \begin{equation}\label{property-of-n2}
    N/e^{nt}<\gamma, \ \ \ \forall n\geq n_2.\end{equation}
    We now use the fact that $(Y,S,\tau)$ satisfies  \Cref{def:good}. This property ensures that we can choose $n_3$ such that for every $n\geq n_3$, the proportion of words $w\in L_{n,s}(Y)$ with $r_n(w)\geq n_2$ is at least $1/\ell$. We claim that \begin{equation}\label{bound-A_n/a_n}\frac{A_n(\epsilon)}{a_n(t)}<\gamma(\ell+2), \ \ \ \forall n\geq n_3.\end{equation}
    Recall that $A_n(\epsilon)$ is equal to the sum of $\spa(T,\tau^{\{0,\dots,n-1\}}(w),\epsilon)$, with $w$ ranging over $L_{n,s}(Y)$. We will separate this sum in the cases $r_n(w)<n_1$ and $r_n(w)\geq n_1$. If $r_n(w)\geq n_1$ then by \Cref{eq:estimate-1} we have  
    $\spa(T,\tau^{\{0,\dots,n\}}(w),\epsilon)\leq \gamma e^{q_n(w)t}$. Therefore \begin{align*}\sum_{\substack{w\in L_{n,s}(Y)\\ r_n(w)\geq n_1}}\spa(T,\tau^{\{0,\dots,n-1\}}(w),\epsilon)&\leq \gamma\sum_{\substack{w\in L_{n,s}(Y)\\ r_n(w)\geq n_1}} e^{q_n(w)t} \\ &\leq \gamma a_n(t).
    \end{align*}
    We now consider the  words $w$ with $r_n(w)< n_1$.  By definition of $N$ in this case we have $\spa(T,\tau^{\{0,\dots,n-1\}}(w),\epsilon)\leq N$. Taking the sum over words $w$ with $r_n(w)< n_1$ we have
    \[\sum_{\substack{w\in L_{n,s}(Y)\\ r_n(w)< n_1}}\spa(T,\tau^{\{0,\dots,n-1\}}(w),\epsilon)\leq |L_{n,s}(Y)|N.
    \]
    Since the proportion of words in $L_{n,s}(Y)$ with $r_n(w)\geq n_2$ is at least $1/\ell$, we have $|L_{n,s}(Y)|\leq |\{w\in L_{n,s}(Y) : r_n(w)\geq n_2\}|\cdot (\ell+1)$. Therefore the inequality above can be continued as
    \begin{align*}
    &\leq |\{w\in L_{n,s}(Y) : q_n(w)\geq n_2\}| (\ell+1) N\\&=(\ell+1)\sum_{\substack{w\in L_{n,s}(Y)\\ r_n(w)\geq  n_2}} N
    \end{align*}
    Our choice of $n_2$ ensures that for every $w$ with $r_n(w)\geq n_2$ we have $N<\gamma e^{q_n(w)t}$ (by \Cref{property-of-n2} and because  $r_n(w)\leq q_n(w)$). Therefore we can bound the previous sum as 
    \(
    \leq (\ell+1) \sum_{\substack{w\in L_{n,s}(Y)\\ r_n(w)\geq  n_2}} \gamma e^{q_n(w)t}
    \leq (\ell+1)\gamma a_n(t)
    \)
    We have proved that for all $n\geq n_3$ we have  $A_n(\epsilon)\leq (\ell+2) \gamma a_n(t)$ and therefore $A_n(\epsilon)/a_n(t)\leq (\ell+2) \gamma$. This shows that the limsup of $A_n(t)/a_n(t)$ as $n\to\infty$ is at most $(\ell+2) \gamma$. Since $\gamma$ is arbitrary, it follows that this limsup us $0$. Thus we have proved \Cref{eq:limsup}. Since this holds for  $\epsilon$ arbitrarily small, we have proved our claim that $\upperent_{\textbf{a}}(S\rtimes_\tau T)\leq h_{top}(T)$. 
    We will now prove that $h_{top}(T)\leq \lowerent_{\textbf{a}}(T)$. This is clearly true if $h_{top}(T)=0$, so we assume now $h_{top}(T)\in (0,\infty]$. We fix an arbitrary $t$ with $0<t<h_{top}(T)$ and we prove that $t\leq \lowerent_{\textbf{a}}(S\rtimes_\tau T)$. According to the definition of lower slow entropy we must prove that for arbitrarily small $\epsilon$ we have
    \[\liminf_{n\to\infty} \frac{\spa(S\rtimes_\tau T,\{0,\dots,n-1\},\epsilon)}{a_n(t)}>0.\]
    Thanks to \Cref{prop:the-set-A}, it suffices to prove that for arbitrarily small $\epsilon$ we have \begin{equation}\label{liminf}\liminf_{n\to\infty} \frac{A_n(\epsilon)}{a_n(t)}>0.\end{equation}
    Choose $t'$ such that $t<t'<h_{top}(T)$. By \Cref{estimate}, for $\epsilon$ arbitrarily small we have that
    \begin{equation}\label{lower-bound-spa}
    \spa(T,F,\epsilon)\geq e^{C(F)|F|t'}, \ \ \ F\Subset \Z.
    \end{equation}
    We will prove that for this $\epsilon$, \Cref{liminf} holds. For each $n\in\N$ we define $g(n)$ as 
    \[g(n)=\min \{\frac{\spa(T,\tau^{\{0,\dots,n-1\}}(w),\epsilon)}{e^{q_n(w)t}} : w\in L_{n,s}(Y)\}\]
    We also let $w_n$ be a word realizing this minimum, that is \[g(n)=\frac{\spa(T,\tau^{\{0,\dots,n-1\}}(w_n),\epsilon)}{e^{q(w_n)t}}.\]
    A direct computation shows that  $g(n)\cdot a_n(t)\leq A_n(\epsilon)$. Therefore $A_n(\epsilon)/a_n(t)\geq g(n)$ for all $n$, and to prove \Cref{liminf} it suffices to show
    \[\liminf_{n\to\infty} g(n)>0.\]
    By \Cref{lower-bound-spa},  for each $n$ we have $\spa(T,\tau^{\{0,\dots,n-1\}}(w_n),\epsilon)\geq e^{q(w_n)t'}$, so
\(\liminf_{n\to\infty} g(n)=\liminf_{n\to\infty} \frac{\spa(T,\tau^{\{0,\dots,n-1\}}(w_n),\epsilon)}{e^{q(w_n)t}}\geq \liminf_{n\to\infty} \frac{e^{q(w_n)t'}}{e^{q(w_n)t}}
=\liminf_{n\to\infty}e^{q(w_n)(t'-t)}
\)
Since $q(w_n)\geq 1$ for each $n$ and $t-t'>0$, this liminf is nonzero.  Thus we have proved \Cref{liminf}. 


Finally, a general fact is that $\lowerent_{\textbf{a}}(S\rtimes_\tau T)\leq \upperent_{\textbf{a}}(S\rtimes_\tau T)$. Thus we have
\[h_{top}(T)\leq \lowerent_{\textbf{a}}(S\rtimes_\tau T)\leq \upperent_{\textbf{a}}(S\rtimes_\tau T)\leq h_{top}(T).\]
This proves that $\lowerent_{\textbf{a}}(S\rtimes_\tau T)$ and $\upperent_{\textbf{a}}(S\rtimes_\tau T)$ are both equal to $h_{top}(T)$.
\end{proof}

\end{document}
